

\begin{abstract}



% Personalized image generation involves creating images tailored to an individual’s visual preferences.  The ability to do so, however, remains inaccessible to the vast majority of users because existing methods require dense user histories or explicit feedback, limiting access to narrow creative demographics. 
% It is well known that visual preferences can vary dramatically between users such that an image generated with the same prompt may delight some while disappointing others.  This is in part because current diffusion models optimize for average aesthetic appeal across a broad audience, rather than individual taste.
% %

% In this paper, theorize and demonstrate that concise descriptors of user's persona can serve as powerful proxies for eliciting visual preferences without any user-specific data. To this end, we introduce a new paradigm, \zipp (Zero-shot Image Personalization from Personas), that communicates users preferences through natural language persona—rich, structured textual descriptions.  \zipp employs a large vision-language model to rewrite user prompts based on the users persona, then uses these prompts to guide generative models to generate personalized content, without weight updates or extra supervision. In doing so, \zipp enables broad generalization to new users, produces interpretable and editable prompts, and works seamlessly across diffusion models while preserving output fidelity.

% To support our new paradigm, we introduce a Zero-Shot Image Personalization Benchmark \zipbench, containing 1.5K users and 40K images, to systematically evaluate persona-driven personalization. Experimental results show that few-shot persona role-play with a Multimodal Large Language Model (MLLM) matches or exceeds state-of-the-art fine-tuned baselines on the BLIFT Visual Behavior Simulation benchmark. Our trained MLLM achieves strong generalization on preference datasets, such as Pick-a-Pic, HPD-v2, and PIP.
% %
% Our results demonstrates that natural-language personas enable universal, scalable, zero-shot personalization, opening generative vision to the global user base without sacrificing interpretability or control.

Personalized image generation adapts diffusion models to individual visual preferences, yet remains inaccessible to most users because existing methods require dense interaction histories or explicit feedback. Current models optimize for average aesthetic appeal rather than individual taste, such that the same prompt may delight some users while disappointing others. We introduce \zipp (Zero-shot Image Personalization from Personas), demonstrating that natural language personas—concise descriptors of user identity, interests, and aesthetic sensibilities—serve as powerful proxies for visual preferences without requiring user-specific data. \zipp employs a large vision-language model to rewrite prompts based on personas, guiding diffusion models to generate personalized content without weight updates or supervision. This enables generalization to unseen users, produces interpretable prompts, and works seamlessly across diffusion architectures.

To support this paradigm, we release \zipbench, containing 1.5K users with mined personas and 40K generated images. Few-shot persona role-play with MLLMs matches fine-tuned baselines on BLIFT Visual Behavior Simulation, while our trained MLLM generalizes strongly to Pick-a-Pic, HPD-v2, and PIP. Our results demonstrate that natural-language personas enable universal, scalable, zero-shot personalization without sacrificing interpretability or control.

\end{abstract}



% Generative models have made remarkable progress in producing photorealistic and semantically aligned images from text. Yet, they remain largely impersonal—optimized for broad aesthetic appeal rather than for the diverse preferences of individual users. Existing personalization methods rely on subject-driven fine-tuning or limited user feedback from small, homogeneous groups, leaving open the question of how to generalize personalization to real-world, demographically diverse users.

% We introduce Zero-Shot Image Personalization, a new paradigm where an MLLM elicits user visual preferences directly from descriptive natural language personas, enabling personalization without any prior user data or model fine-tuning. We further propose a graph neural network (GNN)-based behavioral clustering framework that transforms large-scale posting and commenting behavior into coherent natural-language personas and their corresponding visual preference profiles. Using this method, we construct the largest real-world dataset of user visual preferences, spanning 50K users and 300K images, and release it publicly for research.

% To systematically evaluate this setting, we present the first Zero-Shot Image Personalization Benchmark, which contains 1.5 K users, their personas, and 40 K images, and compare open-source and proprietary LLMs in persona-driven Image personalization. Our MLLM achieves state-of-the-art few-shot performance on the BLIFT Visual Behavior Simulation benchmark and surpasses fine-tuned baselines on existing datasets such as Pick-a-Pic, HPD-v2, and PIP—demonstrating that natural-language personas serve as powerful, scalable proxies for human preference modeling. Our work opens a new direction toward universal, zero-shot personalization in generative vision.